\section{Secondary experiments and failed instruments}
\label{app:secondary}

This appendix explains experiments that validate the apparatus, diagnose a main-track result, test a failed measurement instrument, or provide a reduced external reproduction.
Each subsection states the question, intervention or comparison, estimand and inference unit, result, conclusion, limitation, and role in the thesis.

\subsection{E001: end-to-end implementation validation on a planted synthetic signal}
\label{sec:e001}

E001 asked whether the initial implementation could recover a known relationship from synthetic language features and synthetic fMRI-like targets.
The run exercised feature extraction, ridge fitting, cross-validation, and result serialization while holding the planted data-generating relationship under experimenter control.
Its estimand was recovery of that constructed relationship, and the relevant unit was the synthetic sample rather than a participant or brain measurement.
The implementation recovered the planted signal and therefore passed its engineering validation. \evd{E001}
This result licenses use of the implementation as a starting point, but it provides no evidence that a language model predicts real neural activity, that the nuisance controls are sufficient on real stimuli, or that a neural objective can improve training.
E002 is the first scientific evidence because it replaces the constructed target with recorded language-network responses and adds trained-versus-untrained and nuisance controls.

\subsection{E010 and E010b: participant averaging as a target intervention}
\label{sec:e010}

E010 asked whether the apparent brain-specific intervention contrast increased as more participants were averaged into the target.
The motivation was the discrepancy between E005, which used a five-participant average, and E008, which found a near-zero mean across individual participants.
For nested averages of $k\in\{1,2,3,5,9\}$ participants, the measured intervention gaps were approximately $-0.0002,-0.0001,+0.0023,+0.0194,+0.0071$. \evd{E010}
The original interpretation treated the rising $k=1$-to-$5$ limb as a reliability dose response while noting the reversal at $k=9$.
That interpretation was not identified because every larger target contained the smaller target, participant order was fixed, and target construction changed together with signal-to-noise ratio.

E010b therefore repeated the analysis over random participant subsets at each $k$.
The resulting subset distributions overlapped substantially and did not support a stable monotone causal relation between participant count and the intervention contrast. \evd{E010}
The estimand in both analyses was the contrast for a constructed group-average endpoint, while participant subsets, not seed--fold cells, supplied the sensitivity variation.
The joint conclusion is that averaging changes both reliability and the biological target, and the nested curve does not establish that increasing $k$ causes a brain-specific training effect.
Its role is diagnostic: it explains why E005 and E008 can differ without treating either endpoint as universally correct.

\subsection{E014: averaging improves encoding signal without establishing an intervention artifact}
\label{sec:e014}

E014 separated the effect of participant averaging on the encoding measurement from its effect on brain-guided training.
It held the model features and encoding procedure fixed while replacing individual targets with averages over increasing numbers of participants.
Raw encoding unique $R^2$ rose across $k=1,2,3,5,9$ as $+0.0039,+0.0162,+0.0324,+0.0444,+0.0702$. \evd{E014}
This is the expected consequence of suppressing participant-specific and measurement noise in a shared stimulus-evoked target.
The estimand is nevertheless different at each $k$: an individual response at $k=1$ and an increasingly group-shared response as $k$ grows.
E014 therefore supports participant averaging as a legitimate way to improve the signal-to-noise ratio of a group-level encoding target.
It does not show that averaging universally manufactures a false encoding score, and it does not rescue participant-level inference for E005.
Its role is to keep the SNR benefit separate from the biological-generalization question answered by E008.

\subsection{E016 prerequisite gates and the synthetic-to-real transfer test}
\label{sec:e016}

E016 asked whether a dense stimulus-predictable neural proxy could overcome the scarcity of measured fMRI targets and provide a usable auxiliary training signal.
The program used TRIBE v2 to generate a high-dimensional synthetic neural target from text-associated stimuli, trained language-model students with that target, and then evaluated whether the resulting improvement transferred to an independent real-brain endpoint.

The first prerequisite was fidelity: a synthetic target is useful only if it captures stable stimulus-dependent structure and can be predicted by the selected student representation.
The fidelity analyses showed that the proxy was learnable enough to support a controlled training experiment, but a separate attempt to use it as a formal ceiling on real-brain alignment failed because proxy predictions captured only a limited and target-dependent portion of recorded response variance. \evd{E016}
The ceiling analysis is therefore a failed instrument rather than evidence that no additional stimulus-predictable neural structure exists.

The training comparison used six stochastic seeds and the same KD backbone across target arms.
On their respective synthetic endpoints, the TRIBE-guided arm improved target $R^2$ over KD by \result{e016_tribe_synth_gain}, whereas the sentence-local projected teacher-feature control improved by \result{e016_textfeat_synth_gain}. \evd{E016}
E016 established only output-dimension matching for the text-feature target; it did not itself establish information, geometry, baseline-difficulty, or effective-rank matching.
The block-permuted control also retained some fixed points and is described as block-permuted rather than fully deranged.

The independent transfer gate scored saved students on the fixed Tuckute target averaged over five participants and five functional ROIs.
The inference unit was the six stochastic training seeds for this one endpoint, not participants.
Direct TRIBE-minus-KD was \result{e016_direct_transfer}, with five of six seed estimates negative \evd{E016}.
The relative TRIBE-minus-text-feature gain was \result{e016_transfer_delta}; all six margins were negative and the exact sign test gave $p=0.03125$ \evd{E016}.
E026 later shows why that relative value cannot identify target content.
The corresponding comparison against permuted-control gains was also negative, and PCA ranks 25, 50, and 100 did not reverse the route.
Synthetic training used layer 6 while the Tuckute diagnostic read layer 7, which further bounds mechanistic interpretation.
The conclusion is a transfer failure in this fixed diagnostic, not a population-level null and not a universal upper bound on synthetic neural supervision.
E016 has a specific role in the thesis: learnability of a neural-looking target does not establish usefulness for an independently measured biological endpoint.

\begin{figure}[H]
  \centering
  \safeimage[width=0.94\textwidth]{../figures/fig08_synthetic_transfer.png}
  \caption{The frozen E016 seed-level diagnostic. Panel A shows target-$R^2$ gains over KD on each arm's own synthetic endpoint. Those values are not a common content scale. Panel B shows the relative TRIBE-minus-text-feature change on one fixed participant-averaged Tuckute endpoint. All six values are negative, but the text-feature comparator later fails the E026 measured-axis audit. E025 supplies the separate participant-level result \evd{E016}\evd{E025}\evd{E026}.}
  \label{fig:synthetic-transfer}
\end{figure}

\subsection{E026: retrospective target-control identification audit}
\label{sec:e026}

E026 audits the existing E016 targets and students without retraining and without using Tuckute participant outcomes \evd{E026}.
Two disjoint deterministic 4,096-item samples provide fixed geometry checks, while the six existing seeds quantify only variability of the frozen training procedure.
The mechanical conjunction has 55 required measured-axis checks: 30 fail, 24 pass, and the lag-1024 order check is unresolved because there are no valid within-document pairs.
Mean KD-only target $R^2$ is \result{e026_kd_targets} for TRIBE and textfeat, respectively, leaving a TRIBE-to-textfeat remaining-headroom ratio of \result{e026_headroom_ratio}.
Normalized spectrum log-RMS differences are \result{e026_spectrum_rms} across the two samples; TRIBE versus textfeat stable ranks are descriptively \result{e026_stable_ranks}, with absolute log-ratios \result{e026_stable_rank_log_gaps}; and low-level nuisance-predictability gaps are \result{e026_lowlevel_gaps}. The spectrum differences, stable-rank log-ratios, and nuisance gaps each exceed their frozen engineering margins.
All six global parameter-update ratios and all six layer-6 centered-Frobenius movement ratios also fail, whereas four of six centered-CKA distance checks pass.
Standardized global scale, most valid temporal-order checks, binned power shape, and static-unique nuisance predictability pass.
Initial-gradient comparability is unknown because the gradients were not retained.
The result shows that the saved comparator fails the predeclared measured-axis conjunction; it does not prove information inequivalence, assign causality to any mismatch, establish brain-specific TRIBE content, or explain E025's near-zero direct TRIBE-minus-KD contrast.
Consequently, the larger cross-target synthetic gain cannot carry a specificity claim, while the within-TRIBE learning gain remains valid.

\subsection{E019: external Negi-style objective under local controls}
\label{sec:e019}

E019 asked whether a faithful-as-feasible implementation of a published Negi-style brain-tuning head would first reproduce a local gain and then survive matched-quality and target-specificity controls.
The implementation combined the available naturalistic LeBel substrate with a differentiable temporal head and tested learning-rate, probe, and positive-control variants before the final comparison.
The relevant estimand was held-out alignment change for the real target relative to a permuted target under the same training and evaluation path.
The head did not produce a stable target-specific improvement on the local substrate, so the planned matched-quality analysis of a reproduced positive could not be performed. \evd{E019}
Across three participants and three training seeds, gentle tuning changed encoding correlation by $-0.0006\pm0.0027$ and unique $R^2$ by $-0.0001\pm0.0010$ relative to the base model, where the terms after $\pm$ are descriptive standard deviations across the nine participant--seed cells. \evd{E019}
The real-minus-permuted specificity contrast was $-0.0004\pm0.0018$, and the real-minus-intended-quality-control contrast was $-0.0010\pm0.0018$ on the same descriptive cell scale. \evd{E019}
The intended quality control did not achieve a match: the real-target arm had perplexity about $51.1$--$51.7$, while the control was about $40.1$--$40.3$, so the latter contrast is not a matched-quality estimate. \evd{E019}
Participants, not the nine participant--seed cells, are the biological inference units; these values are mechanism checks rather than population confidence intervals.
Stronger learning rates produced negative encoding changes, and the strongest probe increased perplexity from approximately $50$ to $11{,}514$ while reducing encoding correlation by about $0.055$. \evd{E019}
This corroborates the mechanism failure seen in the thesis's own objectives because a distinct objective family also failed to establish movement.
It does not reproduce the published dataset, participant pool, full training scale, or exact implementation, and therefore cannot refute Negi et al.
Its role is convergent boundary evidence: the local null is not confined to one MSE loss, but external validity remains unresolved.

\subsection{E020: an empirical stimulus-predictable ceiling that did not close}
\label{sec:e020}

E020 asked whether an empirical reference constructed from other participants could approximate $\mathbb{E}[Y\mid S]$ and reveal brain response structure beyond the stimulus-predictable mean.
The intended comparison residualized a target participant's response against a leave-one-participant-out group reference, then measured how much language-model signal remained.
An initial residual appeared to preserve substantial alignment, which would have suggested a participant-specific component beyond the shared stimulus response.
Diagnostics showed that the reference was itself unreliable ($0.33$), residual split-half reliability was $0.174$, and temporal leakage and autocorrelation could preserve apparent signal after subtraction. \evd{E020}
Once the cleanest controlled construction was used, the trained-minus-untrained residual gap was $-0.018$. \evd{E020}
The inference used six participants on one fixed story, and the reference did not provide a high-reliability estimate of the latent conditional mean.
E020 therefore rejects this particular ceiling instrument while leaving the general ceiling question open.
Its role is to prevent the Discussion from turning a failed measurement into either a positive participant-specific mechanism or a universal closure result.

\subsection{E021: reading-time content versus statistical shape}
\label{sec:e021}

E021 asked whether human reading time supplies a cognition-specific auxiliary target after removing the component predictable from language-model surprisal.
The early comparisons used raw reading time, surprisal-residualized reading time, permuted targets, random targets, word length, frequency, and other-model surprisal, and evaluated frozen-probe learning curves across training seeds.
Raw and surprisal-residualized reading time performed similarly, showing that removal of the model's own surprisal-predictable component did not remove the apparent benefit. \evd{E021}
In the penultimate analysis the residual target beat several learnability-matched controls, but those controls did not match its strong temporal autocorrelation and heavy-tailed marginal distribution.
The final control phase-randomized the signal while preserving its autocorrelation and approximate marginal shape but removed behavioral content.
That content-free shape twin tied the reading-time residual: the difference in area under the frozen-probe learning curve was $+5.3$ with a 95\% interval $[-7.9,+18.5]$, where lower area means faster or better learning. \evd{E021}
The independent units were stochastic training seeds, and the claim is restricted to this learning-curve setup and target construction.
The conclusion is that the observed advantage is attributable to signal shape or regularization, not demonstrated cognition-specific information.
E021's role is methodological: auxiliary targets require controls matched for learnability and statistical geometry, not only random or permuted labels.

\subsection{E022: reduced Moussa-style speech reproduction}
\label{sec:e022}

E022 asked whether brain-guided tuning of a pretrained speech representation reproduced a reported gain in phoneme classification before investing in a full matched-quality shrink test.
The reduced setup used one LeBel participant, 24 stories, a wav2vec2-base initialization for both arms, and three training seeds.
The primary estimand was the brain-tuned minus pretrained phoneme-F1 difference on held-out speech.
The observed gain was $+0.52$ F1 points with a 95\% interval $[-0.36,+1.40]$, below the predeclared $+2$-point continuation gate and including zero. \evd{E022}
The planned matched-twin shrink test was consequently vacuous and was not built.
Because this setup is smaller and differs from the published multisubject experiment, it is a reduced non-reproduction rather than evidence that the published effect is false.
Its role is to show that the thesis did not obtain an external positive on which to demonstrate a matched-quality collapse.

\subsection{E024: privileged gaze information stopped at the precondition gate}
\label{sec:e024}

E024 asked whether gaze or EEG observed only during training could improve the sample efficiency of a text classifier, following the learning-using-privileged-information setting.
The planned estimand was the area between held-out learning curves for a text-only student and privileged-information arms at paired training subsets.
Before building five training arms, the design predeclared gaze as the first binding gate and required it to be reliable, informative about the label under natural reading, and evaluable with a leakage-free grouped split.

The reliable gaze summaries largely tracked sentence length and did not predict the relation label under natural reading.
Across binary and multiclass probes, linear and nonlinear readouts, participant-averaged and within-participant features, natural-reading performance stayed near chance; the rich multiclass probe gave balanced accuracy $0.202$ against chance $0.200$ with permutation $p=0.35$. \evd{E024}
In task-directed reading, by contrast, gaze predicted relation type with balanced accuracy $0.295$ against chance $0.100$ and permutation $p=0.003$. \evd{E024}
This contrast indicates task-intent leakage rather than a portable natural-reading signal.

The dataset also contained 300 sentences drawn from only seven paragraphs, so a paragraph-disjoint test set varied from 8 to 161 sentences across split seeds and did not support a stable low-data learning-curve estimand.
These two failures were independent: the signal was uninformative for the clean target, and the grouping structure prevented a stable leakage-free evaluation.
The five-arm LUPI intervention was therefore stopped at its binding gaze precondition gate, and EEG was not evaluated.
This is not a completed negative intervention or an EEG precondition result; it is evidence that the proposed gaze-first substrate cannot identify the intended transferable privileged-information effect.

\subsection{E023: objective-specific shedding as an unexecuted identification program}
\label{sec:e023-detail}

E023 formalized a counterfactual question left open by E003 and E015: whether a distilled model lies below the alignment expected for a normally trained model of the same language-model quality.
It proposed a same-lineage quality frontier, a matched-quality KD-versus-ordinary-fine-tuning comparison, and operation in a regime where the frontier is sufficiently flat to separate objective effects from ordinary quality loss.
The only executed component was an analysis-only slope and knee gate using already measured Pythia and Qwen models.
The Pythia 1B-to-410M pair met the local flatness screen with slope $+0.015$, alignment change $+0.0007$, and quality change $0.066$ bits per byte. \evd{E023}
The three checked Qwen pair slopes were $+0.334$, $-0.217$, and $+0.076$, so none met the predeclared saturated-regime requirement \evd{E023}.
The analysis-only gate marked a lower-quality pilot as feasible but did not clear the planned training comparison; no inferential unit or hypothesis test was claimed for this descriptive screen.
No KD-versus-fine-tuning intervention was run, so E023 contributes an identification framework rather than a scientific result about objective-specific shedding.
